\section{Abelian sheaves}
\label{section-abelian-sheaves}

\noindent
Let $\mathcal{C}$ be a site.
The category of abelian sheaves on $\mathcal{C}$ is denoted
$\textit{Ab}(\mathcal{C})$. It is the full subcategory of
$\textit{PAb}(\mathcal{C})$ consisting of those abelian presheaves
whose underlying presheaves of sets are sheaves.
Properties ($\alpha$) -- ($\zeta$) of
Sites, Section \ref{sites-section-sheaves-algebraic-structures}
hold, see
Sites,
Proposition \ref{sites-proposition-functoriality-algebraic-structures-topoi}.
In particular the inclusion functor
$\textit{Ab}(\mathcal{C}) \to \textit{PAb}(\mathcal{C})$
has a left adjoint, namely the sheafification functor
$\mathcal{G} \mapsto \mathcal{G}^\#$.

\medskip\noindent
We suggest the reader prove the lemma on a piece of scratch paper rather
than reading the proof.

\begin{lemma}
\label{lemma-abelian-abelian}
Let $\mathcal{C}$ be a site. Let $\varphi : \mathcal{F} \to \mathcal{G}$
be a morphism of abelian sheaves on $\mathcal{C}$.
\begin{enumerate}
\item The category $\textit{Ab}(\mathcal{C})$ is an abelian category.
\item The kernel $\Ker(\varphi)$ of $\varphi$ is the same as the
kernel of $\varphi$ as a morphism of presheaves.
\item The morphism $\varphi$ is injective
(Homology, Definition \ref{homology-definition-injective-surjective})
if and only if $\varphi$ is injective as a map of presheaves
(Sites, Definition \ref{sites-definition-presheaves-injective-surjective}),
if and only if $\varphi$ is injective as a map of sheaves
(Sites, Definition \ref{sites-definition-sheaves-injective-surjective}).
\item The cokernel $\Coker(\varphi)$ of $\varphi$ is the sheafification
of the cokernel of $\varphi$ as a morphism of presheaves.
\item The morphism $\varphi$ is surjective
(Homology, Definition \ref{homology-definition-injective-surjective})
if and only if $\varphi$ is surjective as a map of sheaves
(Sites, Definition \ref{sites-definition-sheaves-injective-surjective}).
\item A complex of abelian sheaves
$$
\mathcal{F} \to \mathcal{G} \to \mathcal{H}
$$
is exact at $\mathcal{G}$ if and only if for all
$U \in \Ob(\mathcal{C})$ and all $s \in \mathcal{G}(U)$
mapping to zero in $\mathcal{H}(U)$ there exists a covering
$\{U_i \to U\}_{i \in I}$ in $\mathcal{C}$ such that each
$s|_{U_i}$ is in the image of $\mathcal{F}(U_i) \to \mathcal{G}(U_i)$.
\end{enumerate}
\end{lemma}

\begin{proof}
We claim that Homology, Lemma \ref{homology-lemma-adjoint-get-abelian}
applies to the categories $\mathcal{A} = \textit{Ab}(\mathcal{C})$
and $\mathcal{B} = \textit{PAb}(\mathcal{C})$, and the functors
$a : \mathcal{A} \to \mathcal{B}$ (inclusion),  and
$b : \mathcal{B} \to \mathcal{A}$ (sheafification).
Let us check the assumptions of
Homology, Lemma \ref{homology-lemma-adjoint-get-abelian}.
Assumption (1) is that $\mathcal{A}$, $\mathcal{B}$ are additive categories,
$a$, $b$ are additive functors, and $a$ is right adjoint to $b$.
The first two statements are clear and adjointness is
Sites, Section \ref{sites-section-sheaves-algebraic-structures} ($\epsilon$).
Assumption (2) says that $\textit{PAb}(\mathcal{C})$ is abelian
which we saw in Section \ref{section-abelian-pre-sheaves} and
that sheafification is left exact, which is
Sites, Section \ref{sites-section-sheaves-algebraic-structures} ($\zeta$).
The final assumption is that $ba \cong \text{id}_\mathcal{A}$ which is
Sites, Section \ref{sites-section-sheaves-algebraic-structures} ($\delta$).
Hence Homology, Lemma \ref{homology-lemma-adjoint-get-abelian}
applies and we conclude that $\textit{Ab}(\mathcal{C})$ is abelian.

\medskip\noindent
In the proof of Homology, Lemma \ref{homology-lemma-adjoint-get-abelian}
it is shown that $\Ker(\varphi)$ and $\Coker(\varphi)$
are equal to the sheafification of the kernel and cokernel of $\varphi$
as a morphism of abelian presheaves. This proves (4). Since the kernel
is a equalizer (i.e., a limit) and since sheafification commutes with
finite limits, we conclude that (2) holds.

\medskip\noindent
Statement (2) implies (3). Statement (4) implies (5) by our description
of sheafification. The characterization of exactness in (6) follows from
(2) and (5), and the fact that the sequence is exact if and only if
$\Im(\mathcal{F} \to \mathcal{G}) =
\Ker(\mathcal{G} \to \mathcal{H})$.
\end{proof}

\noindent
Another way to say part (6) of the lemma is that a
sequence of abelian sheaves
$$
\mathcal{F}_1 \longrightarrow
\mathcal{F}_2 \longrightarrow
\mathcal{F}_3
$$
is exact if and only if the sheafification of
$U \mapsto \Im(\mathcal{F}_1(U) \to \mathcal{F}_2(U))$
is equal to the kernel of $\mathcal{F}_2 \to \mathcal{F}_3$.

\begin{lemma}
\label{lemma-limits-colimits-abelian-sheaves}
Let $\mathcal{C}$ be a site.
\begin{enumerate}
\item All limits and colimits exist in $\textit{Ab}(\mathcal{C})$.
\item Limits are the same as the corresponding limits of abelian presheaves
over $\mathcal{C}$ (i.e., commute with taking sections over objects of
$\mathcal{C}$).
\item Finite direct sums are the same as the corresponding finite direct sums
in the category of abelian pre-sheaves over $\mathcal{C}$.
\item A colimit is the sheafification of the corresponding colimit in
the category of abelian presheaves.
\item Filtered colimits are exact.
\end{enumerate}
\end{lemma}

\begin{proof}
By Lemma \ref{lemma-limits-colimits-abelian-presheaves} limits and colimits
of abelian presheaves exist, and are described by taking limits and colimits
on the level of sections over objects.

\medskip\noindent
Let $\mathcal{I} \to \textit{Ab}(\mathcal{C})$, $i \mapsto \mathcal{F}_i$
be a diagram. Let $\lim_i \mathcal{F}_i$ be the limit of the diagram
as an abelian presheaf. By Sites, Lemma \ref{sites-lemma-limit-sheaf}
this is an abelian sheaf. Then it is quite easy to see that
$\lim_i \mathcal{F}_i$ is the limit of the diagram in
$\textit{Ab}(\mathcal{C})$. This proves limits exist and (2) holds.

\medskip\noindent
By Categories, Lemma \ref{categories-lemma-adjoint-exact}, and because
sheafification is left adjoint to the inclusion functor we see that
$\colim_i \mathcal{F}$ exists and is the sheafification of the colimit
in $\textit{PAb}(\mathcal{C})$. This proves colimits exist and (4) holds.

\medskip\noindent
Finite direct sums are the same thing as finite products in any abelian
category. Hence (3) follows from (2).

\medskip\noindent
Proof of (5). The statement means that given a system
$0 \to \mathcal{F}_i \to \mathcal{G}_i \to \mathcal{H}_i \to 0$
of exact sequences of abelian sheaves over a directed set $I$ the sequence
$0 \to \colim \mathcal{F}_i \to \colim \mathcal{G}_i \to
\colim \mathcal{H}_i \to 0$ is exact as well. A formal argument using
Homology, Lemma \ref{homology-lemma-check-exactness} and the
definition of colimits shows that the sequence
$\colim \mathcal{F}_i \to \colim \mathcal{G}_i \to \colim \mathcal{H}_i \to 0$
is exact. Note that $\colim \mathcal{F}_i \to \colim \mathcal{G}_i$
is the sheafification of the map of presheaf colimits which is
injective as each of the maps $\mathcal{F}_i \to \mathcal{G}_i$ is
injective. Since sheafification is exact we conclude.
\end{proof}







